% =============================================================================
% MOTOR, FUENTE E IDIOMA
% =============================================================================

% Este proyecto debe compilarse con XeLaTeX.
\usepackage{fontspec}
\setmainfont{Times New Roman}

\usepackage[spanish,es-tabla]{babel}
\usepackage{csquotes}
\usepackage{datetime}

% =============================================================================
% DISEÑO DE PÁGINA Y TEXTO
% =============================================================================

\usepackage[
    a4paper,
    top=2.54cm,
    bottom=2.54cm,
    left=2.54cm,
    right=2.54cm
]{geometry}
\usepackage{setspace}
\usepackage[none]{hyphenat}

% =============================================================================
% FIGURAS, TABLAS Y COLOR
% =============================================================================

\usepackage{graphicx}
\usepackage{float}
\usepackage{caption}
\usepackage{tikz}
% La librería babel evita el choque con los caracteres activos del español
% («<» y «>» para las comillas angulares), que de otro modo rompen «>=» en tikzset.
\usetikzlibrary{arrows.meta,positioning,calc,fit,backgrounds,shapes.geometric,shapes.misc,decorations.pathreplacing,babel}

\usepackage{array}
\usepackage{booktabs}
\usepackage{longtable}
\usepackage{tabularx}
\usepackage[table]{xcolor}

\definecolor{lightgray}{gray}{0.9}

% =============================================================================
% MATEMÁTICAS
% =============================================================================

\usepackage{amsmath}
\usepackage{amssymb}

% =============================================================================
% ÍNDICES, TÍTULOS Y ENCABEZADOS
% =============================================================================

\usepackage{tocloft}
\usepackage{fancyhdr}
\usepackage{titlesec}

% =============================================================================
% URL, ENLACES Y BIBLIOGRAFÍA
% =============================================================================

% xurl amplía los puntos de corte de url y carga el paquete url.
\usepackage{xurl}
\urlstyle{same}
\Urlmuskip=0mu\relax

\usepackage[
    style=ieee,
    backend=biber,
    sorting=none,
    maxbibnames=99,
    minbibnames=1,
    giveninits=true,
    doi=true,
    url=true,
    isbn=false,
    eprint=false
]{biblatex}

% Formato bibliográfico requerido por el documento.
\DeclareFieldFormat[article,inproceedings,misc,techreport]{title}{\textit{#1}}
\DeclareFieldFormat{url}{\url{#1}}
\DeclareFieldFormat{doi}{%
    \ifhyperref
        {\href{https://doi.org/#1}{\nolinkurl{https://doi.org/#1}}}
        {\nolinkurl{https://doi.org/#1}}%
}
\DeclareNameAlias{author}{family-given}

\setlength{\bibitemsep}{0.5\baselineskip}
\setlength{\bibhang}{0pt}

\addbibresource{bibliography/references.bib}

% hyperref se carga al final para que pueda adaptar los comandos anteriores.
\usepackage[hidelinks,breaklinks=true]{hyperref}
